% Fixture: real TeX glyph positions, scripts, fraction rules, limits and matrix delimiters.
% Rebuild with tectonic formula-layout.tex. This source is the transcription reference.
\documentclass[11pt]{article}
\usepackage{amsmath,amssymb}
\usepackage[margin=1in]{geometry}
\pagestyle{empty}
\begin{document}
\section*{Formula transcription fixture}
The formulas below are separated by ordinary prose so their reading regions can be identified.
\begin{equation}
\hat\beta=(X^TX+\lambda I)^{-1}X^Ty
\end{equation}
The residual sum of squares has indexed terms and limits on a displayed sum.
\begin{equation}
\mathrm{RSS}(\beta)=\sum_{i=1}^{N}(y_i-x_i^T\beta)^2
\end{equation}
The derivative below contains a fraction rule and a superscript inside an operator's argument.
\begin{equation}
\frac{\partial}{\partial x}\sin(x^2)=2x\cos(x^2)
\end{equation}
The entries of a two by two matrix must retain their row and column positions.
\begin{equation}
A=\begin{pmatrix}a&b\\c&d\end{pmatrix}
\end{equation}
The next expression places limits above and below an integral.
\begin{equation}
\int_0^\infty e^{-x}\,dx=1
\end{equation}
The radical's rule covers both terms of its argument.
\begin{equation}
r=\sqrt{x^2+y^2}
\end{equation}
The final expression contains a fraction nested inside another fraction.
\begin{equation}
f(x)=\frac{1}{1+\frac{1}{x}}
\end{equation}
\end{document}
